\section{Conclusion}
\label{sec:conclusion}
% Other members' paragraphs belong alongside this contribution.
\paragraph{Calibration and inspection.}
Forecast confidence changes the price of a routing decision even when the spam label stays
fixed. In the example in Section~\ref{sec:calibration}, moving the higher forecast above
the spam-folder threshold raises cost from $0.25$ to $1.10$ per email. In the Spambase
experiment, inspecting roughly one in four emails at price $0.3$ reduced realized loss for
both fitted models. The appropriate comparison depends on the cost of an error, the price of
review, and how much probability mass actually enters the inspection band.

\paragraph{Forecast quality and the prior correction.}
% Member 5
A cost-sensitive filter acts on the value of a forecast, not only on its side of $0.5$. On Spambase,
rescaling the logistic model's log-odds left its accuracy at $93.6\%$ but raised its cost by a third, and only
the probability scores noticed. Calibration should be checked before deployment; Platt scaling on validation
data is a cheap repair. When only the spam rate changes, one constant shift of the log-odds restores the
Bayes rule (Section~\ref{sec:t5}).

\paragraph{Accuracy is the wrong target, and inspection has a price ceiling.}
% Member 6
In the analytic model of Section~\ref{sec:t6} every quantity is exact, so the effects are not noise. The most
accurate rule ($84.6\%$) costs $0.772$ per email, while the cost-optimal rule is less accurate ($73.9\%$) and
costs $0.311$: when the two errors differ in price, accuracy and cost rank rules in opposite order. Inspection
helps exactly when it is cheaper than the worst automatic loss, $\ci<\cstar=\cfp\cfn/(\cfp+\cfn)=10/11$; at
$\ci=0.3$ it cuts the cost from $0.311$ to $0.179$. Every experiment makes and scores its decisions with the
same shared code, and the threshold rule matches a brute-force minimisation on $10{,}000$ random cases, so
differences between models are not artefacts of how decisions are made or measured.
